\subsection{A carrier that survives the low-degree cleanup of one resolved cell}

\begin{theorem}[AT24: Proved; a single-cell higher-arity cleanup with actual maps]
\label{AT24}
Work over $R=\mathbb Z_{(2)}$ and $S=\Gamma_R(T)$.
Let $X\to A,B$ be actual strict maps with the source and
target hypotheses of AT20. Suppose the marked cell is
an actual free C21 cell in normalized simplex degree two:
its generators $x,y=dx$ have forms $1,2$, weight $b$,
and boundary $\partial x=dh$.
In particular its simplex-two algebra has the actual
presentation
\[
 X_2=C\otimes_R\Lambda_R(x)\otimes_R\Gamma_R(y),
 \qquad dx=y,
 \tag{AT.157}
\]
where, as a cochain algebra, $C$ is $S$ tensored over
$R$ with the free PD algebra on the other identity-disk
generators over $R$.
The differential on $C$ is the one in that presentation.
Use the bounds of AT20, put $s=2b$, and form the
torsionfree resolved $I_1\langle s\rangle$ extension $Y$
of AT16 for $(\gamma_2(y),xy)$.
Choose its target maps by the actual AT20 construction.

There is an actual sequence of strict inclusions, made
from two-cell $K(2)$ extensions, torsionfree images and
the based saturation A-SAT01, whose union $Y^\sharp$
has extending strict maps to both $A$ and $B$.
It has the following properties.
\begin{enumerate}
\item Its terms are torsionfree, its rational algebra is
      a coefficient resolution in every simplex and
      weight, and its coefficient $H^0$ is $S$.
\item Its simplex-zero and simplex-one algebras are
      literally those of $X$.
\item In simplex two its coefficient comparison is
      an isomorphism in forms zero through three,
      in every weight.
\item In simplex two no generator of cell arity zero,
      one or two is added after the original resolved extension.
      Here the selected $x,y$ have cell arity one,
      every original resolved generator has arity two,
      and the other generators in $C$ have arity zero.
\item For every target there remains an actual strict
      PD cochain map from the simplex-two algebra to
      the target's simplex-three algebra, extending
      the explicitly constructed carrier below.
\end{enumerate}
All added two-cell repairs kill actual torsion classes
of forms two or three and cell arity at least three.
The support $q\leq w$ is retained.

If the initial $X$ has horizontally acyclic rows in
the eventual ranges of C21, the same horizontal
zero- and one-homology ranges are retained throughout.
The C21 construction can then restore all horizontal
rows using only cells of simplicial dimension at least
three, with extending target maps; it leaves
$Y^\sharp_2$ unchanged.
This last construction does not assert vertical
exactness in forms two or three in higher simplices.

The construction treats one selected free odd cell.
It does not assert the analogous arity property for
products of distinct old cells, or after arbitrary
later degree-four repairs.
\end{theorem}

\begin{proof}
We first construct the carrier on the whole
simplex-two algebra. Fix one of the targets, denoted
$A$. Use the actual elements $q,\alpha,x',y'$ in
AT20, and write
\[
 \widetilde x=f_A(x)-dq,\qquad
 a_3=s_0x'+\alpha,\qquad
 d_0a_3=\widetilde x,\quad d_1a_3=x',
 \quad d_2a_3=d_3a_3=0.
\]
The actual normalized $D(3)$ coefficient-resolution
lift constructed there gives values for every
resolved generator in $A_3$.
On the algebra $C$ in (AT.157) use $s_0f_A$.
On the two selected old generators use
\[
 x\longmapsto a_3+d(s_0q),\qquad
 y\longmapsto da_3.
 \tag{AT.158}
\]
On $v_0(t)$ subtract
$\varepsilon(t)(s_0q)\,da_3$ from the normalized
value supplied by the auxiliary twist; leave all
other resolved generator values as constructed there.
The same augmentation equation
$\varepsilon\partial=0$ used in AT20 verifies every
coefficient relation, including the terminal
$W^0$ relation. The differential of the correction is
exactly the change from $\widetilde x\,y$ to $xy$.
Thus (AT.158) and these values extend to an actual
strict PD cochain map
\[
                  F_A:Y_2\longrightarrow A_3.
 \tag{AT.159}
\]
The freeness in (AT.157) is used here: no assignment
to a marked element of an arbitrary algebra is
silently extended across relations.
The torsionfree image causes no difficulty because
$A_3$ is torsionfree.

Its zero face is the chosen $f_{A,2}$.
Its first face on the indicated generators is
\[
 \begin{array}{c|c}
  \text{generator}&d_1F_A\\ \hline
  x&x'+dq\\
  y&y'\\
  v_0(t)&-\varepsilon(t)q\,y'\\
  u_0(t),\ u_i(t),\ v_i(t)\ (i>0),\ W^0&0 .
 \end{array}
 \tag{AT.160}
\]
On $C$ this face is $f_A$.
The second and third faces kill every positive-arity
generator in this table.
These identities follow directly from
$d_1s_0q=q$, $d_2s_0q=d_3s_0q=0$ and the
normalization of all uncorrected resolved values.
They include the actual old generator $x$, rather
than only $\widetilde x$.

At simplex two the differential and all defining
relations are homogeneous for cell arity.
The torsionfree image is homogeneous too, since
it is the image in the corresponding rationally
graded algebra. In (AT.160) each possibly nonzero
positive-arity value has form degree at least its
assigned arity: the pairs are $(1,1),(2,1),(2,2)$.
The remaining values vanish.
Consequently
\[
 d_1F_A(D^{a,r})=0\quad(r<a),\qquad
 d_2F_A(D^{a,r})=d_3F_A(D^{a,r})=0\quad(a>0)
 \tag{AT.161}
\]
for $D=Y_2$. Multiplication, divided powers and sums
preserve this conclusion. The degrees of coefficients
from $C$ are nonnegative and their arity is zero.

Here is the actual repair step preserving (AT.161).
Let $q_*=2$ or $3$, let $a\geq3$, and suppose
$z\in D^{a,q_*}$ is closed and
$b_0\in D^{a,q_*-1}$ satisfies
\[
                         db_0=mz,\qquad m\ne0.
 \tag{AT.162}
\]
Every element of cell arity at least two has all
simplex-two faces zero. Indeed $y$ and all original
resolved generators have zero faces; the only old
positive-arity generator that may have a nonzero
face is the single odd $x$, whose square is zero.
All subsequently added generators below have zero
faces as well. Thus $z,b_0$ define actual linear
$K(2)$ maps by their normalized values.

By (AT.161), $F_A(b_0)$ has all three positive faces
zero, because $q_*-1<a$.
Its differential is $mF_A(z)$.
Torsionfreeness then proves that $F_A(z)$ also has
these faces zero. Hence both belong to $N_3(A)$,
with zero faces their specified images under $f_A$.
Their weight is at least $ab\geq3b$.
The bounds of AT20 give
$3b\geq s+1\geq\lambda_A(4)$.
The exact normalization functor therefore gives
\[
             H^j(N_3(A)^{w,*})=0\qquad(0\leq j\leq4)
             \quad\text{for }w\geq s+1.
 \tag{AT.163}
\]
Here positive vertical cohomology vanishes by
admissibility, and degree zero vanishes because
$N_3$ of the constant coefficient module is zero.

Choose actual integral normalized elements
$U_A,V_A\in N_3(A)$ with
\[
                  dU_A=F_A(z),\qquad
                  dV_A=mU_A-F_A(b_0).
 \tag{AT.164}
\]
Their form degrees are $q_*-1,q_*-2$.
The first equation uses (AT.163) in degree $q_*$,
and the second in degree $q_*-1$.
The two zero faces $d_0U_A,d_0V_A$ have all their
faces zero by the simplicial identities.
They therefore give the actual target values for
the $K(2)$ two-cell extension
\[
       du=z,\qquad dv=mu-b_0,\qquad
       |u|=q_*-1,\quad |v|=q_*-2,
       \quad\operatorname{arity}(u)=\operatorname{arity}(v)=a.
 \tag{AT.165}
\]
Repeat independently in $B$, on the same source.
At simplex two extend the carrier by
$F_A(u)=U_A,F_A(v)=V_A$.
All their positive faces vanish, so (AT.161)
is retained. This is an actual extension by the
free PD universal property.

Rationally put $u'=u-b_0/m$. Then $du'=0$ and
$dv=mu'$. This is an actual invertible triangular
change of the two free generators.
The free rational graded-commutative disk has its
usual monomial contraction. Thus the rational
coefficient comparison is retained, in every simplex.
Taking the image in the rational algebra retains
the old torsionfree source injectively and all
target maps. New weights are at least $3b$.
The target's coefficient $H^0$ in that range
therefore proves the integral coefficient $H^0$,
as in AT14. The parameters $K(2)$ vanish in
simplices zero and one.
All these assertions apply to an arbitrary set of
simultaneous extensions (AT.165), since each element
uses finitely many generators.

We also check the first skeleton during saturation.
The free identity-disk presentation of $X_2$ gives
$H^1(X_2)=0$. Since $d_0s_0=1$, $X_1$ is an actual
cochain retract of $X_2$ and has $H^1(X_1)=0$.
The coefficient $H^0(X_1)$ is $S$ by hypothesis.
The $K(2)$ steps leave $X_1$ literally fixed.
For a rational simplex-one root, its integral
differential therefore has a primitive in $X_1$;
adjusting its coefficient to the prescribed vertex
value makes that primitive the original root.
Thus the proof of A-SAT01 adds no simplex-one
element at any stage. This uses only the stated
$H^0$ and $H^1$ facts, not unproved higher exactness
of $X_1$.

We next check saturation, including its compatibility
with the carrier. At simplex two the original arity
zero part is $C$. Its positive cohomology in degrees
one through three vanishes: an even-to-odd free
identity disk is integrally contractible, while an
odd-to-even disk has arity-$n$ differential
$R\xrightarrow{\,n\,}R$ in degrees
$n(r+1)-1,n(r+1)$ when its odd generator has degree
$r\geq1$. Its first positive cohomology is in degree
at least four. Tensoring these actual two-term
complexes proves the same assertion for $C$.
There are no new arity-one generators, and that
part is the identity disk $x\mapsto y$ tensored
with $C$.

In arity two the old selected disk has the
two-term complex
$Sxy\xrightarrow{\,2\,}S\gamma_2(y)$ in degrees
three and four.
The resolved linear complex has sole cohomology
$I_1\langle s\rangle/2$ in degree three.
Its twisting map on cohomology is the actual
inclusion
\[
             I_1\langle s\rangle/2
                    \hookrightarrow S\langle s\rangle/2 .
 \tag{AT.166}
\]
The long exact sequence of these two explicitly
given complexes proves vanishing in degrees zero
through three for their combined arity-two part.
Tensoring with $C$ preserves this vanishing.
For clarity this last assertion is over the PID
$R$: $C$ is $S$ tensored with the free $R$ complex
on its other disk generators, whose $H^0$ is $R$
and whose positive cohomology begins in degree four.
The combined arity-two complex also has first
cohomology in degree four and has no $H^0$.
In total degrees at most three both the tensor
and the $\operatorname{Tor}_1$ terms vanish.
This elementary two-term free-presentation
calculation is also valid for infinite generator
sets, by filtered unions of the indicated complexes.

It follows that saturation adds no arity-zero,
arity-one or arity-two element in simplex two.
For example, an arity-two rational zero-form with
integral differential has an integral primitive
by the just proved $H^1$ vanishing; their rational
difference is zero because this arity has no $H^0$.
The other two arities are treated by the preceding
coefficient and identity-disk calculations.
Every positive-arity zero-form has zero vertex
values, since its monomials contain a positive-arity
zero-form generator, all of whose faces are zero.
Here is why saturation performed in all simplices
still preserves this simplex-two grading.
Write a rational zero-form root in simplex two as
the finite sum $t=\sum_a t_a$ of its arity components.
The differential preserves arity and the current
simplex-two algebra is arity-graded, so each $dt_a$
is integral. Every positive-arity zero-form has all
faces zero, by its positive-arity zero-form factors;
the same holds rationally. Thus each $t_a$, $a>0$,
has based vertex value zero, while $t_0$ has the
integral based value of $t$ itself.
Every $t_a$ is therefore one of the roots explicitly
adjoined by (SAT.1). Conversely their sum is such
a root. The PD closure is consequently arity-graded.
Faces of roots adjoined in higher simplices are
already roots of (SAT.1) in simplex two, by the
proved simplicial stability in A-SAT01; they add no
exception to this decomposition.
Its new components all have arity at least three,
by the preceding low-arity calculation.
Every such component still has weight at least $ab$.

Let $t$ be one of the new homogeneous rational
zero-form roots of arity $a\geq3$.
The rational carrier $F_{A,\mathbb Q}(t)$ has
integral differential. By $H^1(A_3)=0$ in its
weight it has an integral primitive for that
differential. Its vertex-zero value is zero, so
the coefficient adjustment in (SAT.3) proves
that $F_{A,\mathbb Q}(t)$ itself is integral.
This extends the carrier through every PD
closure and every saturation stage.
Moreover $dt$ has form one and arity $a>1$,
so (AT.161) makes the differential of every
positive face of $F_{A,\mathbb Q}(t)$ zero.
Those faces are rational coefficients, with
zero vertex value, hence are zero.
Thus the new zero-form roots have zero reference
faces, and (AT.161) is retained.
The equality $d_0F_A=f_{A,2}$ is retained because
both extensions are restrictions of the previous
rational maps. The argument applies independently
to both targets and uses the fixed cutoff $s+1$
of A-SAT01.

Now saturate $Y$, and iterate the following
operation countably many times.
For every nonzero homogeneous cohomology class
of forms two or three in simplex two choose
a representative $z$, a nonzero annihilator $m$,
and a homogeneous $b_0$ as in (AT.162); attach
(AT.165) and saturate again.
Rational coefficient comparison supplies these
annihilators, and the low-arity calculation shows
that every such class has arity at least three.
It also places its weight at least $3b$, so the
same target bounds suffice at every stage.
The actual maps and carriers constructed above
are compatible with all inclusions.

Take the filtered union of these stages.
Every closed form of degree two or three lies
in one stage. If its class there is nonzero,
the following repair step kills it; if it is
already a boundary its primitive remains.
Every closed one-form is killed by the next
saturation stage. This proves the full claimed
low-degree comparison in simplex two.
The coefficient, rational, torsionfree, support
and first-skeleton assertions hold at every
stage and hence in the union.
No arity at most two changes.

Finally all the inclusions fix simplices zero
and one and retain the old simplex-two elements.
Thus row zero-homology is unchanged, and row
one-homology is a quotient of the old one.
Their stated eventual vanishing ranges are
therefore retained. To restore the remaining
rows, run the actual C21 construction only on
normalized cycles of degree at least two,
using a cutoff in each form range at least
the old two-row cutoff and both target row
cutoffs. The target horizontal exactness
provides its actual generator values.
Every added cell then has normalized parameter
dimension at least three and does not change
simplex two. The established C21 argument proves
the asserted row exactness after its iteration.
Its generators are chosen only in weights at least
their form degrees, exactly as in C21, retaining
the support condition.
In each simplex its new free generator complements
are vertical identity disks; their rational monomial
contractions retain the rational coefficient comparison.
This uses precisely the relative free-disk construction
of C21 and does not posit a new lifting statement.
This last step supplies no new assertion about
higher-simplex vertical cohomology.
\end{proof}
